\section{General complex exclusion of \texorpdfstring{$\sha{R8_3_3472}$}{R8-3-3472}}
\label{sec:r3472-focus}
The ring $R=\sha{R8_3_3472}$ is excluded by the complete
central-induction and rotation argument below. Its necessary domains
and exact terminal certificate are presented in one proof sequence.

\begin{theorem}\label{thm:r3472-general-exclusion}
The fusion ring $\sha{R8_3_3472}$ has no fusion categorification over
$\mathbb C$. In particular, the conclusion does not require a pivotal,
spherical, braided, unitary, or pseudo-unitary structure as an initial
hypothesis.
\end{theorem}
\begin{proof}
Theorem~\ref{thm:research-four-general-spherical} reduces a hypothetical
complex categorification to a pseudo-unitary Galois conjugate with the
same fusion rules. We recall that reduction, then exclude every possible
center restriction frame. There are $372$ labelled low frames; exact
span and phase certificates leave three equivalent frames. Their twists
belong to a complete domain of $424$ Galois orbit types. The necessary
rotation equations for these types have an exact integer contradiction,
proved by $1648$ rational bound deductions and a final separating
inequality. The details and replay scope follow.
\end{proof}

\subsection{Reduction from arbitrary complex categorification}
\label{subsec:r3472-exclusion-scope}
All eight source simples are selfdual. In the canonical sphericalization
of an arbitrary categorification, both lifts of each simple are selfdual,
and changing the lift reverses its second indicator
\cite[Section~5.1]{OstrikPiv3}. Choose the lift of indicator $+1$.
The odd-multiplicity theorem~\cite[Theorem~3.9]{LPRV} makes the
multiplicity of every negative lift in a positive-lift square even.
Consequently the signed quotient tensor satisfies
\[
 |T_{ij}^k|\leq N_{ij}^k,\qquad
 T_{ij}^k\equiv N_{ij}^k\pmod2,\qquad
 T_{jj}^i\equiv N_{jj}^i\pmod4.
\]
The complete search, with the unit, duality, cyclic trace symmetry,
dual reversal and all ordered associativity equations, has $84$ signed
variables and $298865973$ nodes. Its only solution is $T=N$.
No further sign gauge or commutativity assumption is added. The positive
lifts therefore form a spherical section, proving that the original
category admits a spherical structure.

Of the eight ordinary characters of $R$, the six outside the FP Galois
orbit have zero values on simple basis elements. They cannot be spherical
dimension characters. The remaining pair is the FP character and its
$\sqrt2$-conjugate. Extend the required embedding to an automorphism of
$\mathbb C$ and transport the associator, rigidity and spherical structure
simultaneously. All coherence identities are preserved, and the resulting
spherical dimensions are positive FP dimensions. Thus the transported
category is pseudo-unitary. This also avoids any separate assumption that
a chosen pivotal structure descends to a number field.
The original gate, source hashes and complete replay remain in
\path{evidence/rank8_research_20261022/pivotality/}; the new proof package
contains an independent audit and fresh successful replay of this step.

\subsection{The complete distinguished restriction frames}
Put
\[
 r=\sqrt2,\quad \delta=7+5r,\quad
 \pi=\delta^2=99+70r,\quad \beta=1+\pi=100+70r,
 \quad \gamma=7\delta=49+35r,\quad D=7\beta.
\]
The source dimensions, in canonical order, are
$(1,\delta-1,\delta-1,\delta,\delta-1,\delta+1,\delta+1,\delta+1)$.
Let $V=X_1+\cdots+X_7$. We use $F:Z(\mathcal C)\to\mathcal C$ and
its adjoint $I$, and write $v_Z$ for the nonnegative integral restriction
column of a central simple $Z$.
The source formal codegrees are $D,D'=700-490r$, and six copies of $7$.
Central induction gives eight distinct twist-one summands $A_E$ of
$I(\one)$, with dimensions $1,\pi,\beta,\ldots,\beta$
\cite[Theorem~2.13]{OstrikPiv3}.

The restriction of the distinguished object of dimension $\pi$ is forced:
\begin{equation}\label{eq:r3472-fixed-pi}
 F(A_{D'})=\one+2V.
\end{equation}
To avoid an identification assumption, first let $P$ be the image of the
vacuum under a modular Galois automorphism that flips $r$. The Galois
vacuum row gives $\dim(P)=D/D'=\pi$ and, by the induction character
formula~\cite[Lemma~4.3 and Corollary~4.5]{DeshpandeCenter},
\[
 \chi_E(F(P))=\begin{cases}
 \pi&E=\mathrm{FP},\\
 \pi^{-1}&E=\mathrm{FP}',\\
 1&f_E=7.
 \end{cases}
\]
The element $\one+2V$ has exactly these eight character values.
Character inversion proves $F(P)=\one+2V$. Its unit multiplicity is one,
so $P$ belongs to $I(\one)$ and is the distinguished $A_{D'}$.
This proves \eqref{eq:r3472-fixed-pi} without presuming that no other
central object has dimension $\pi$.

The induction Gram matrix is the tensor-derived integer matrix
$B=\sum_iN_iN_i^{\mathsf T}=\sum_Zv_Zv_Z^{\mathsf T}$.
Set $u=e_0$, $p=(1,2,\ldots,2)^{\mathsf T}$ and
$Q=B-uu^{\mathsf T}-pp^{\mathsf T}$. A distinguished $\beta$ column
$v$ has $v_0=1$ and $\sum_{i>0}v_i=14$. With
$A=v_1+v_2+v_4$ and $C=v_5+v_6+v_7$, its dimension equation gives
$C=A+1$ and $v_3=13-2A$, where $0\leq A\leq6$.
These equations give exactly $2142$ possible columns. The necessary
condition $Q-vv^{\mathsf T}\succeq0$ retains $96$.
For the six labelled columns, impose distinguished $S$-symmetry, the
pairwise Gram inequalities, their prescribed sum $Q_{0,*}$, and positivity
of the residual Gram matrix. Exact exhaustive enumeration leaves $372$
labelled frames. Repeated restriction columns are allowed throughout.
A separate enumeration and exact positive-semidefinite test reproduce
this complete list.

The high columns, those with $v_0=0$, are reconstructed from the source
rather than taken from a previously feasible witness. The complete
induction bounds, modular eigenvalue integrality and Galois spectral
support tests leave $90$ restrictions of dimension
$\eta=\beta/2$, $78$ of dimension $\beta$, and one each of dimensions
$c\gamma$ for $c=1,2,5,10$; the latter columns are $cV$.
Let $a,b$ count the high $\eta,\beta$ simples, and put
$w=\sum c^2$ over the other high simples. Comparing global dimension
and orthogonality of the dimension character with its $r$-conjugate gives
\[
 50a+200b+49w=8402,\qquad
 50a+200b-49w=-1202.
\]
Hence $a+4b=72$ and $w=98$, excluding $c=10$.
The resulting $171$ high restriction possibilities form a complete
necessary domain. The source-derived domain replay checks each
integrality and missing-Galois-conjugate rejection exactly.

\subsection{Reduction of all frames to one orbit}
Let $X=(\chi_E(X_j))$ be the full character table and let $L$ have
columns $v_{A_E}$ in the same character order. The distinguished
modular $STS$ identities and induction give
\begin{equation}\label{eq:r3472-full-sts-restriction}
 \sum_Z\theta_Zv_Zv_Z^{\mathsf T}=LX.
\end{equation}
For each frame, the high columns must lie in the support of its residual
positive-semidefinite Gram matrix. Exact linear annihilators show that
\eqref{eq:r3472-full-sts-restriction} cannot lie in the span of the
supported rank-one matrices for $276$ frames.
Nine verified based-ring automorphisms partition the remaining $96$
frames into twelve orbits. Eleven exact rational separating certificates
exclude eleven orbits. The only frames left are $42,117,304$, one orbit
under these automorphisms.

These phase certificates do not assume a maximum twist order. Apply
normalized relative trace to
$K=\mathbb Q(r,\zeta_7+\zeta_7^{-1})$ in
\eqref{eq:r3472-full-sts-restriction}. The $33$ exact trace values of
roots whose orders divide $280$, together with their convex hull, contain
the normalized relative trace of every root of unity. Cyclotomic tower
transitivity and the factor $\mu(m)/\varphi(m)$ for the part of the
order coprime to $14$ prove the coverage. Introducing nonnegative real
masses on these atoms is therefore a necessary relaxation for arbitrary
twists. Each separating covector is checked against every supported
column and atom, as well as the exact right-hand side. The frame-coverage
replay verifies all $276$ span certificates, the nine automorphisms,
all eleven phase certificates, and the three surviving labels.

\subsection{Maximal third indicators}
\begin{proposition}\label{prop:r3472-maximal-third-indicator}
In the canonical positive spherical structure of every pseudo-unitary
categorification of $\sha{R8_3_3472}$,
the third Frobenius--Schur indicator of every nonunit simple is $2$.
Equivalently, rotation acts identically on each two-dimensional space
$\operatorname{Hom}(1,X_i^{\otimes3})$, $i>0$.
\end{proposition}
\begin{proof}
Put $\delta=7+5\sqrt2$, $D=700+490\sqrt2$, and $D'=700-490\sqrt2$.
The source dimensions are
$(1,\delta-1,\delta-1,\delta,\delta-1,\delta+1,\delta+1,\delta+1)$,
and $\dim\operatorname{Hom}(1,X_i^{\otimes3})=2$ for all $i>0$.
The spherical Cauchy theorem~\cite[Theorem~3.9]{BNRWCauchy} and
$DD'=9800$ imply that $3$ is coprime to
$N=\operatorname{FSexp}(\mathcal C)$.
The dimensions lie in $\mathbb Q(\zeta_N)$, so $8\mid N$.
For $3\widetilde3\equiv1\pmod N$, the automorphism
$\sigma(\zeta_N)=\zeta_N^{\widetilde3}$ therefore flips $\sqrt2$.
The higher Gauss-sum formula~\cite[Theorem~4.4, (4.17)]{NSWHigherGauss}
gives
\[
 \sum_i d_i\nu_3(X_i)=\frac D{\sigma(D)}
 =\frac D{D'}=\delta^2=1+2\sum_{i>0}d_i.
\]
We have $\nu_3(1)=1$. For $i>0$, $\nu_3(X_i)$ is the trace of an
order-three operator on a two-dimensional space, so
$\operatorname{Re}\nu_3(X_i)\le2$.
All $d_i$ are positive, so the displayed equality forces equality
for every summand. Every rotation eigenvalue is consequently $1$;
the operator is diagonalizable because its cube is the identity.
\end{proof}

\subsection{Rigidity of the surviving frame}
Use frame $42$, and identify high columns with their seven nonunit
coordinates. Its residual Gram matrix has null-vector
$(1,1,-6,1,1,1,1)^{\mathsf T}$. Thus every high column satisfies
$\sum_{i>0}v_i=7v_3$. An $\eta$ column would give $v_3=1$,
$A+C=6$ and $C-A=1$, impossible over the integers. Hence there are
$18$ high $\beta$ simples and $\sum c^2=98$ for the remaining objects.
For a high $\beta$ column, $v_3=2$, $A=5$, and $C=7$.
Writing $z=v-2V$ gives $\|z\|^2\geq2$. The residual Gram trace is
$1226$ and forces $\sum\|z\|^2=36$, so equality holds for all
$18$ columns. Consequently the only possibilities are
\[
 2V-e_a+e_c,\qquad a\in\{1,2,4\},\quad c\in\{5,6,7\}.
\]
The cross entries of the Gram matrix force each of these nine columns to
occur twice. The remaining columns are $cV$, $c\in\{1,2,5\}$,
with dimensions $c\gamma$.

The ten associated rank-one matrices are linearly independent.
Solving \eqref{eq:r3472-full-sts-restriction} makes the twist sum on
each two-object $\beta$ column equal to $2\cos(2\pi k/7)$, where,
in the displayed row and column orders, the $k$-matrix is
\[
 \begin{pmatrix}1&3&2\\2&1&3\\3&2&1\end{pmatrix}.
\]
A pair of unit complex numbers with this sum is
$\{\zeta_7^k,\zeta_7^{-k}\}$. Thus the $18$ high $\beta$ twists
are fixed, while the six distinguished $\beta$ twists are one.
For the other simples define
\[
 P_n=\sum_{\dim Z=c\gamma}c^2\theta_Z^n.
\]
The same $STS$ computation and the second and third indicators give
\begin{equation}\label{eq:r3472-first-three-moments}
 P_0=98,\qquad P_1=-2-10r,\qquad P_2=-2,\qquad P_3=-2+10r.
\end{equation}
For the second indicators, odd-square witnesses for objects
$1,2,3,4,5,6,7$ are respectively squares of objects
$4,1,1,2,6,7,5$; Theorem~3.9 of~\cite{LPRV} gives $\nu_2=1$.
The third indicators are all two by
Proposition~\ref{prop:r3472-maximal-third-indicator}.
The central indicator formula~\cite{NgIndicators} now becomes, for
$i>0$,
\begin{equation}\label{eq:r3472-indicator-moments}
 \nu_n(X_i)=b_i(n)+\frac{r}{20}(P_n+2),\qquad
 b_i(n)=\begin{cases}
 1&7\nmid n,\\
 6&i\in\{1,2,4\},\ 7\mid n,\\
 7&i=3,\ 7\mid n,\\
 8&i\in\{5,6,7\},\ 7\mid n.
 \end{cases}
\end{equation}

\subsection{A complete finite twist domain}
Cauchy's theorem~\cite[Theorem~3.9]{BNRWCauchy} and
$DD'=9800$ restrict exponent primes to $2,5,7$.
Modular Galois symmetry sends a twist $\theta$ to $\theta^{a^2}$
for every unit $a$ modulo the exponent. It preserves each dimension
$c\gamma$: under the $r$-flip, $D/\sigma(D)=\pi$ and
$\sigma(\gamma)=-\gamma/\pi$, with the Galois sign correcting the
minus sign. Therefore the weighted multiplicities in $P_n$ are constant
on unit-square orbits.

For an order $q=2^a5^b7^c$, the orbit size is the product of
\[
 s_2(a)=2^{\max(a-3,0)},\qquad
 s_5(b)=\begin{cases}1&b=0,\\2\cdot5^{b-1}&b>0,\end{cases}
 \quad
 s_7(c)=\begin{cases}1&c=0,\\3\cdot7^{c-1}&c>0.\end{cases}
\]
Since the entire weighted mass is $98$, every occurring orbit has size
at most $98$. It follows that $a\leq9$, $b\leq3$, $c\leq2$, with
the product bound still imposed. There are exactly $63$ possible orders
and $424$ unit-square orbit types. This is a proved finite domain, not a
conductor cutoff. For each type $\mathcal O$ introduce its nonnegative
integer weighted multiplicity $x_{\mathcal O}$; a categorical realization
would have $x_{\mathcal O}=\sum_c c^2m_{c,\mathcal O}$. Allowing every
nonnegative integer $x_{\mathcal O}$ only weakens the necessary system.
Then
\[
 P_n=\sum_{\mathcal O}x_{\mathcal O}
             \sum_{\theta\in\mathcal O}\theta^n,
 \qquad \sum_{\mathcal O}|\mathcal O|x_{\mathcal O}=98.
\]
The power sums are evaluated by exact local Gauss periods. They lie in
$\mathbb Q(\zeta_8,\sqrt5,\sqrt{-7})$; computation in
$\mathbb Q(\zeta_{1120})$ is only a common coordinate representation,
not a restriction that twists have order dividing $1120$.
Independent reduction modulo cyclotomic polynomials checks the local
formulas and all $424$ orbit types.

\subsection{Rotation constraints and the exact terminal contradiction}
Let $E_n$ be pivotal cyclic rotation on
$\operatorname{Hom}(\one,X_i^{\otimes n})$. It has order dividing $n$.
For $1\leq k<n$, set $g=\gcd(n,k)$ and $m=n/g$. The rotation identity
\cite[Theorem~5.1 and its proof]{NgHigherPivotal} gives
\[
 \operatorname{Tr}(E_n^k)
   =\sigma_{k/g}\bigl(\nu_m(X_i^{\otimes g})\bigr),
 \qquad \sigma_{k/g}(\zeta_m)=\zeta_m^{k/g}.
\]
The tensor determines $X_i^{\otimes g}$ exactly, and the indicator is
additive. Thus \eqref{eq:r3472-indicator-moments} computes each trace.
The Fourier multiplicities
\[
 \frac1n\sum_{k=0}^{n-1}\zeta_n^{-tk}\operatorname{Tr}(E_n^k)
 \quad(0\leq t<n)
\]
are nonnegative integers. Impose the necessary integrality congruences
from their constant coefficient in the chosen cyclotomic power basis,
for $i=1,3,5$ and
$n\in\{4,5,7,8,10,16,20\}$. Together with
\eqref{eq:r3472-first-three-moments}, this gives $362$ integer equations
in the $424$ orbit variables and $148$ integer quotient variables.
All variable bounds follow from the mass $98$ and these congruences.

There is a further exact necessary condition:
$\nu_n(X_i)\in\mathbb Q(\zeta_n)$.
If $a$ is coprime to the exponent, $a\equiv1\pmod n$, and
$\chi(a)$ is the sign of the action on $r$,
\eqref{eq:r3472-indicator-moments} implies
\[
 P_n-\chi(a)P_{na}=2(\chi(a)-1).
\]
Their coordinate equations add $266$ rows, giving a necessary system
\begin{equation}\label{eq:r3472-integer-system}
 A x=b,\qquad \ell\leq x\leq u,\qquad x\in\mathbb Z^{572},
 \qquad A\in\mathbb Z^{628\times572}.
\end{equation}
Every coefficient was independently reconstructed with arbitrary-precision
arithmetic from the tensor, exact power sums and rotation identities.

The terminal proof uses only rational arithmetic and integer rounding.
For any rational row covector $y$ and coordinate vector $e_j$, the
identity $Ax=b$ gives the valid lower bound
\[
 x_j\geq yb+
 \sum_k\min\bigl((e_j-yA)_k\ell_k,(e_j-yA)_ku_k\bigr).
\]
Rounding its right-hand side upward is valid because $x_j$ is an integer;
applying the same argument to $-x_j$ gives upper bounds.
The certificate records $380+486+782=1648$ such sequential deductions.
For the resulting box $[\ell^*,u^*]$, its final rational covector $y$ satisfies
\[
 yb-\sum_j\max\bigl((yA)_j\ell_j^*,(yA)_ju_j^*\bigr)
 =\frac{2717541156573712541191977062303}
        {359103652930042901916159999600}>0.
\]
This contradicts $Ax=b$. It excludes frame $42$ and its two equivalents,
so no low frame can be the center of a pseudo-unitary categorification.
The general complex reduction proves
Theorem~\ref{thm:r3472-general-exclusion}.

\subsection{Proof record and replay}
The canonical proof package is
\begin{quote}\small\path{evidence/r3472_exclusion_20261022/}.\end{quote}
It binds the canonical source tensor, complete domain and frame audits,
source-derived rotation system, rational cuts and final separator.
The complete replay command is
\begin{quote}\small
\texttt{python3 }\path{evidence/r3472_exclusion_20261022/replay.py}.
\end{quote}
For the self-contained terminal arithmetic certificate alone, run
\begin{quote}\small
\texttt{python3 }\path{evidence/r3472_exclusion_20261022/terminal/check.py}.
\end{quote}
The terminal checker uses Python's standard library and exact fractions;
no optimizer or floating-point feasibility verdict is used as proof.
An independently implemented checker replays the same contradiction.
A separate exact SCIP certificate is also accepted by the unmodified
VIPR checker, with its input independently matched to the audited system.
These arithmetic checks certify the supplied system; the source derivation,
complete frame coverage and general sphericalization reduction are distinct
parts of the full proof and are all included in the record.

Together with Theorem~\ref{thm:r1768-general-exclusion}, the exact-M3 totals are
$3508$ general complex exclusions, $16$ unitary models and three unresolved
rings: $\sha{R8_3_928}$, $\sha{R8_3_941}$ and $\sha{R8_3_1154}$.
Cumulatively through multiplicity three this gives $4404$ exclusions,
$93$ models and three unresolved rings. Necessary-data witnesses and conditional module calculations are not
used as categorical existence claims.

